\section{Conclusion}
\label{sec:conclusion}

Adherence to a national inpatient hyperglycemia guideline can be measured
automatically, recommendation by recommendation, across a hospital network.
Across 48,354 admissions in six Israeli general hospitals, adherence improved
modestly and unevenly after guideline publication. Higher adherence was
associated with lower 30-day mortality and fewer hyperglycemic complications
after balancing measured baseline differences. The least-implemented
recommendations were largely those whose adherence tracked subsequent
outcomes: repeated insulin decisions throughout hospitalization rather than
discrete admission actions. These findings support evaluating fuller
implementation while recognizing the associated increase in hypoglycemia.
They describe associations rather than treatment effects, and guideline
adoption rather than individual clinician performance. Automated assessment
makes these questions answerable across a health system rather than only in
samples small enough for manual review.
